\question{\textbf{(امتیازی)}}

\begin{ابجد}
\item اندازه‌ی فیلتر $2\times2$ هست، پس $F_w = F_h = 2$. ابتدا ابعاد خروجی را محاسبه می‌کنیم:
\begin{align*}
     & \text{\lr{output width}} = \frac{W - F_w + 2P}{S_w} + 1 = \frac{4 - 2 + 2(1)}{1} + 1 = 5  \\
     & \text{\lr{output height}} = \frac{H - F_h + 2P}{S_h} + 1 = \frac{4 - 2 + 2(1)}{1} + 1 = 5
\end{align*}
با افزودن \lr{padding} ها، نقشه‌ی ورودی به این شکل خواهد بود:
\begin{align*}
    X_{\text{padded}} = \begin{bmatrix}
                            \textcolor{red}{0} & \textcolor{red}{0} & \textcolor{red}{0} & \textcolor{red}{0} & \textcolor{red}{0} & \textcolor{red}{0} \\
                            \textcolor{red}{0} & 1                  & 2                  & 0                  & 1                  & \textcolor{red}{0} \\
                            \textcolor{red}{0} & 3                  & 1                  & 2                  & 1                  & \textcolor{red}{0} \\
                            \textcolor{red}{0} & 0                  & 1                  & 3                  & 1                  & \textcolor{red}{0} \\
                            \textcolor{red}{0} & 2                  & 3                  & 1                  & 2                  & \textcolor{red}{0} \\
                            \textcolor{red}{0} & \textcolor{red}{0} & \textcolor{red}{0} & \textcolor{red}{0} & \textcolor{red}{0} & \textcolor{red}{0}
                        \end{bmatrix}
\end{align*}
هر درایه از نقشه‌ی خروجی نیز به این شکل محاسبه می‌شود:
\begin{align*}
    O_{1,1} = \begin{bmatrix} 0 & 0 \\ 0 & 1 \end{bmatrix} \odot \begin{bmatrix} 1 & -1 \\ 0 & 1 \end{bmatrix} = (0 \times 1) + (0 \times -1) + (0 \times 0) + (1 \times 1) = 1
\end{align*}
در نهایت اگر عملیات کانولوشن را به طور کامل انجام دهیم، نقشه‌ی خروجی به این شکل خواهد بود:
\begin{align*}
    Y = \begin{bmatrix}
            1  & 2  & 0  & 1  & 0 \\
            2  & 0  & 4  & 0  & 1 \\
            -3 & 3  & 2  & 2  & 1 \\
            2  & 2  & -1 & 4  & 1 \\
            -2 & -1 & 2  & -1 & 2
        \end{bmatrix}
\end{align*}

\item
\begin{enumerate}
    \item برای محاسبه‌ی ماتریس توجه، باید ماتریس $Q$ را در ترانهاده‌ی $K$ ضرب کنیم و سپس بر اساس بعد بردارها، آن را \lr{scale} کنیم:
          % \item برای محاسبه‌ی ماتریس توجه، باید ماتریس $Q$ را در ترانهاده‌ی $K$ ضرب کنیم (از آنجایی که در صورت سوال ذکر شده مکانیزم توجه از نوع \lr{Dot-Product} می‌باشد، لازم نیست این ماتریس را \lr{scale} کنیم):
          \begin{align*}
              \frac{QK^T}{\sqrt{d_k}} = (\frac{1}{\sqrt{2}}) \begin{bmatrix} 1 & 2 \\ 0 & 1\end{bmatrix} \begin{bmatrix} 0 & 1 \\ 1 & 1\end{bmatrix} = \frac{1}{\sqrt{2}} \begin{bmatrix} 2 & 3 \\ 1 & 1\end{bmatrix} = \begin{bmatrix} \sqrt{2} & \frac{3\sqrt{2}}{2} \\ \frac{\sqrt{2}}{2} & \frac{\sqrt{2}}{2} \end{bmatrix}
              %   {QK^T} = \begin{bmatrix} 1 & 2 \\ 0 & 1\end{bmatrix} \begin{bmatrix} 0 & 1 \\ 1 & 1\end{bmatrix} = \begin{bmatrix} 2 & 3 \\ 1 & 1\end{bmatrix}
          \end{align*}

    \item با اعمال \lr{Softmax} روی حاصل بخش قبل داریم:
          \begin{align*}
              \text{\lr{Softmax}}(\frac{QK^T}{\sqrt{d_k}}) = \begin{bmatrix} \frac{e^{\sqrt{2}}}{e^{\sqrt{2}} + e^{\frac{3\sqrt{2}}{2}}} & \frac{e^{\frac{3\sqrt{2}}{2}}}{e^{\sqrt{2}} + e^{\frac{3\sqrt{2}}{2}}} \\ \frac{e^{\frac{\sqrt{2}}{2}}}{e^{\frac{\sqrt{2}}{2}} + e^{\frac{\sqrt{2}}{2}}} & \frac{e^{\frac{\sqrt{2}}{2}}}{e^{\frac{\sqrt{2}}{2}} + e^{\frac{\sqrt{2}}{2}}} \end{bmatrix} = \begin{bmatrix} 0.33 & 0.67 \\ 0.5 & 0.5 \end{bmatrix}
          \end{align*}

    \item در نهایت برای محاسبه‌ی خروجی نهایی، ماتریس حاصل را در $V$ ضرب می‌کنیم:
          \begin{align*}
              Y = (\text{\lr{Softmax}}(\frac{QK^T}{\sqrt{d_k}})) \times V = \begin{bmatrix} 0.33 & 0.67 \\ 0.5 & 0.5 \end{bmatrix} \begin{bmatrix} 2 & 1 \\ 1 & 3 \end{bmatrix} = \begin{bmatrix} 1.33 & 2.34 \\ 1.5 & 2 \end{bmatrix}
          \end{align*}

\end{enumerate}


\end{ابجد}